\documentclass[10pt,a4paper]{article}
\usepackage[T1]{fontenc}
\usepackage[utf8]{inputenc}
\usepackage[margin=15mm]{geometry}
\usepackage{lmodern,booktabs,amsmath,microtype,titlesec}
\usepackage[hidelinks]{hyperref}
\usepackage{listings}
\usepackage{enumitem}
\setlength{\parindent}{0pt}
\setlength{\parskip}{3pt}
\titlespacing*{\section}{0pt}{9pt}{5pt}
\setlist{nosep,leftmargin=*}
\lstset{basicstyle=\ttfamily\footnotesize,columns=fullflexible,
  keepspaces=true,breaklines=true,showstringspaces=false,
  aboveskip=4pt,belowskip=4pt}
\title{\vspace{-13mm}Scalable Recommendation with Spark ALS\\
  \large Amazon Reviews 2023: ID2221 Project Report}
\author{Group 0: Riccardo Fragale, Tommaso Tron, Edoardo F.A. de Cal}
\date{}

\begin{document}
\maketitle
\vspace{-9mm}

\section*{Implementation and dataset}
We built a reproducible batch recommendation pipeline using PySpark/Spark SQL,
Spark MLlib explicit-feedback ALS, Parquet, MongoDB, and Streamlit. Separate
stages ingest data, sample users, construct temporal splits, train and evaluate
models, benchmark scaling, and publish recommendations. Some run manifests preserve
configuration, input checksums, seeds, code identity, and artifact lineage.
All completed stages are protected against overwriting and recommendation batches can
resume after interruption.

The dataset is the McAuley Lab's \emph{Amazon Reviews 2023}, restricted to the category
\emph{Movies and TV}\footnote{Dataset and category downloads:
\url{https://amazon-reviews-2023.github.io/}Citation: Hou et al.,
\emph{Bridging Language and Items for Retrieval and Recommendation}, arXiv:2403.03952, 2024.}
as otherwise it would be too long and complex for our machines to compute.
The raw category contains 17.3M reviews. We retain user ID, item ID
(\texttt{parent\_asin}), rating (1--5), and timestamp, plus a separate metadata
projection for item titles. Keeping the latest review per user/item pair removes
170K duplicates, leaving 17.15M interactions, 6.5M users, and 750K
items. Duplicates with equal timestamp retain the highest rating. Typed, validated
interactions are stored in 27 yearly Parquet partitions (1997 to 2023; 756 MB).
Unix timestamps are converted in UTC. Malformed records, missing IDs,
out-of-range ratings, and implausible timestamps are counted by rejection reason in case they are found.
Decompressed JSONL permits
parallel input reads, whereas gzip input is not splittable. Metadata are
deduplicated to one row per item so title joins cannot multiply recommendations.

\section*{Methodology}
Development uses a deterministic, seeded 20\% user sample, preserving each selected
user's full cleaned history. A global chronological split targets 80/10/10 by
interaction count, keeping identical timestamps together. Iterative training-only
5-core filtering removes users and items with fewer than five interactions until
convergence. Held-out data are projected onto the retained training vocabulary;
dense, deterministic integer mappings are persisted for ALS.
Sampling uses a stable hash of the original user ID and seed 2221(selected randomly), making cohorts
reproducible and benchmark fractions nested. Split boundaries are selected from
cumulative timestamp-group counts, so actual proportions may differ slightly
from their targets. Filtering is never fitted on validation or test labels.
For the original full-category training window, 5-core filtering retains
5.9M ratings (43\%).

We compare four systems: interaction-count popularity; a regularized bias-only
model; plain explicit ALS; and bias+ALS, which fits signed rating residuals:
\[
  \widehat r_{ui}=\mu+b_u+b_i+p_u^{\mathsf T}q_i.
\]
Biases are fitted by alternating user/item updates, with regularization 10,
at most 10 passes, and tolerance $10^{-4}$. ALS uses 10 iterations. Validation
NDCG@10 selects among ranks $\{10,25,50\}$ and regularizations
$\{0.03,0.1,0.3\}$ for each ALS variant. Selected settings are rank 50 with
regularization 0.03 (ALS) and 0.1 (bias+ALS). Popularity windows of 30, 90, 180,
365 days and all history are compared on validation; the selected window is 180 days.

After settings are frozen, models are refitted on full-category train+validation
history through 17 November 2020, recomputing the 5-core and mappings. This leaves
6.7 training ratings, 598,310 users, and 175,550 items. The later test window
is untouched by fitting. Ranking uses exact full-catalog scores, excluding seen
items, with deterministic ties. Ratings $\geq4$ define relevance; Precision@10
and binary NDCG@10 are macro-averaged over the same 44,612 eligible users. RMSE
uses the same 162,922 training-known test ratings, with predictions clipped to
$[1,5]$. Popularity's rating predictor is a smoothed item mean, separate from its
ranking counts. Cold-start exclusions and serving fallbacks are reported separately.
Precision@10 divides the number of relevant recommendations by ten even for a
short list; NDCG discounts hits by $1/\log_2(r+1)$ and normalizes by the ideal
ordering of up to ten relevant items. Users without an eligible relevant test
item are excluded and counted. Raw predictions remain available for diagnostics,
while ranking scores are not clipped. User-level bootstrap intervals quantify
ranking uncertainty, and a uniformly random unseen-item list supplies a reference.

Exact ALS ranking broadcasts the item-factor matrix and scores user blocks with
NumPy inside Spark partitions, masking seen items before choosing the top items.
For bias+ALS, item biases are included in the ranking score; adding user bias and
the global mean restores full rating predictions. Unlike latent-only retrieval,
this approach cannot omit a high-bias item merely because its residual score is low.

\newpage
\section*{Recommendation quality}
\begin{center}
\begin{tabular}{lrrrr}
\toprule
System & NDCG@10 & Precision@10 & Users with a hit & RMSE \\
\midrule
Popularity (180 days) & 0.00340 & 0.00093 & 415 & 1.263 \\
Bias-only             & 0.00007 & 0.00004 & 17  & \textbf{1.151} \\
Plain ALS             & 0.00015 & 0.00006 & 27  & 1.448 \\
Bias+ALS              & 0.00003 & 0.00002 & 7   & 1.173 \\
\bottomrule
\end{tabular}
\end{center}
Recent popularity is the strongest ranker; its NDCG@10 bootstrap 95\% interval
is $[0.0030,0.0038]$. Bias-only predicts ratings best, beating the training-global-mean
reference (RMSE 1.355). Explicit ALS does not improve on recent popularity for
ranking or on biases for rating prediction. Absolute ranking accuracy remains low;
the sparse retained graph and substantial cold-start exclusions limit coverage.
The random reference expects only 7.1 users with a top-10 hit, compared with 415
for recent popularity and 27 for ALS. Bias+ALS produces seven, approximately the
random level. All-history popularity also produces only seven, indicating that
recency matters in this temporal evaluation. ALS's training RMSE of 0.148 versus
test RMSE of 1.448 indicates poor generalization for the validation-selected
configuration. Selecting for NDCG does not necessarily optimize rating accuracy.

\section*{Scaling, serving, and limitations}
The benchmark we made uses nested 25/50/100\% user samples with shared
temporal cutoffs and frozen model settings. 
Work times exclude Spark startup and verification counts;
ETL is measured separately and MongoDB publication is excluded.
\begin{center}
\begin{tabular}{lrrr}
\toprule
Work time (seconds) & 25\% & 50\% & 100\% \\
\midrule
Temporal split and filtering & 34.3 & 56.0 & 120.6 \\
ALS training                 & 22.6 & 38.7 & 72.3 \\
Bias+ALS training            & 34.0 & 53.3 & 97.4 \\
Exact ALS top-20             & 27.8 & 70.7 & 277.8 \\
Exact bias+ALS top-20        & 27.6 & 74.5 & 335.2 \\
Total, including other stages & 185.6 & 349.9 & 972.6 \\
\bottomrule
\end{tabular}
\end{center}
Total work grows $5.24\times$ for approximately $4\times$ input. Exact ALS-based
ranking accounts for 63\% of work at full scale and grows with users times catalog
size. Changing 5-core retention also changes the workload. These measurements
demonstrate single-machine input scaling, not multi-machine speedup; there are no
repetitions to estimate timing variability. Cluster deployment is unsupported by
the current configuration validator. Item-factor broadcasting and collection of
per-user evaluation metrics also impose memory limits.
The split grows $3.51\times$ for $4\times$ input; ALS training grows
$3.20\times$ for $4.9\times$ retained training rows. In contrast, the retained
user--item catalog product grows $12.5\times$, explaining why exact ranking grows
much faster than training. Event logs record shuffle, spill, garbage collection,
and task skew to support bottleneck analysis. Each stage starts a new Spark
process, and failed or timed-out stages remain visible in benchmark records.

The serving publication uses popularity with a 180 day window for every user,
chosen because of its superior ranking results. All 6.5M users receive 20
unique unseen items: 600K known users and 5.9M fallback users. Thus the
deployed ranking is not personalized. Bounded batches produce Parquet serving
artifacts; idempotent bulk upserts write versioned MongoDB collections. An active
pointer switches only after count and sampled-document validation. Streamlit
looks up original user IDs and displays recommendations and saved experiment
evidence. Missing titles are explicitly shown (using derived metadata labels or IDs).
Fallback users comprise also 5.2M users removed by the 5-core run. Coverage audits check list length, uniqueness,
rank order, finite scores, and exclusion of previously seen items. All published
lists have length 20, but serving coverage must not be confused with the much
smaller training-known cohort used for model-quality comparisons.
Compared with the proposal, exact custom ranking replaces
\texttt{recommendForUserSubset}, and custom Spark metric calculations replace
MLlib evaluators to expose cohort, coverage, and short-list handling.

\section*{Running the code}
Use Python 3.11, Java 17, and \texttt{uv}; Docker Compose is required for MongoDB.
Install dependencies with \texttt{uv sync --frozen}. Configure the MongoDB
credentials using \texttt{.env.example} and export \texttt{MONGODB\_URI}.
After running the pipeline and publishing recommendations, launch Streamlit with
\texttt{make dashboard CONFIG=configs/real.toml} at \url{http://localhost:8501}.

\textbf{Full instructions are provided in \texttt{README.md}}, including dataset
downloads, environment setup, pipeline commands, experiment reproduction, and
MongoDB configuration.
\end{document}
